
\documentclass{article}
\usepackage{colortbl}
\usepackage{makecell}
\usepackage{multirow}
\usepackage{supertabular}

\begin{document}

\newcounter{utterance}

\centering \large Interaction Transcript for game `bbh', experiment `web\_of\_lies', episode 1 with reflexion{-}Qwen3.5{-}27B.
\vspace{24pt}

{ \footnotesize  \setcounter{utterance}{1}
\setlength{\tabcolsep}{0pt}
\begin{supertabular}{c@{$\;$}|p{.15\linewidth}@{}p{.15\linewidth}p{.15\linewidth}p{.15\linewidth}p{.15\linewidth}p{.15\linewidth}}
    \# & \multicolumn{2}{c}{Player} && \multicolumn{2}{c}{Game Master} \\
    \hline

    \theutterance \stepcounter{utterance}  
    & & & \multicolumn{4}{p{0.6\linewidth}}{
        \cellcolor[rgb]{0.9,0.9,0.9}{
            \makecell[{{p{\linewidth}}}]{
                \texttt{\tiny{[P1$\langle$GM]}}
                \texttt{Evaluate a random boolean function expressed as a word problem.} \\
\\ 
\texttt{Q: Question: Fidel tells the truth. Jerry says Fidel tells the truth. Vina says Jerry tells the truth. Millicent says Vina lies. Raymond says Millicent lies. Does Raymond tell the truth?} \\
\texttt{A: Yes} \\
\\ 
\texttt{Q: Question: Kristian lies. Millie says Kristian lies. Maybelle says Millie tells the truth. Fidel says Maybelle lies. Leda says Fidel lies. Does Leda tell the truth?} \\
\texttt{A: Yes} \\
\\ 
\texttt{Q: Question: Kristian tells the truth. Michaela says Kristian lies. Raymond says Michaela tells the truth. Osvaldo says Raymond tells the truth. Jamey says Osvaldo tells the truth. Does Jamey tell the truth?} \\
\texttt{A: No} \\
\\ 
\texttt{Q: Question: Kristian lies. Sherrie says Kristian lies. Delbert says Sherrie lies. Jerry says Delbert tells the truth. Shalonda says Jerry tells the truth. Does Shalonda tell the truth?} \\
\texttt{A:} \\
            }
        }
    }
     \\ \\

    \theutterance \stepcounter{utterance}  
    & \multicolumn{4}{p{0.6\linewidth}}{
        \cellcolor[rgb]{0.561,0.737,0.561}{
            \makecell[{{p{\linewidth}}}]{
                \texttt{\tiny{[P1$\rangle$GM]}}
                \texttt{Let's break down the logic step{-}by{-}step:} \\
\\ 
\texttt{1.  **Kristian lies**: This is given as a fact. So, Kristian = **False**.} \\
\texttt{2.  **Sherrie says Kristian lies**: Since Kristian actually lies, Sherrie's statement is correct. Therefore, Sherrie = **True**.} \\
\texttt{3.  **Delbert says Sherrie lies**: We know Sherrie tells the truth. Delbert claims she lies, which is incorrect. Therefore, Delbert = **False**.} \\
\texttt{4.  **Jerry says Delbert tells the truth**: We know Delbert lies. Jerry claims Delbert tells the truth, which is incorrect. Therefore, Jerry = **False**.} \\
\texttt{5.  **Shalonda says Jerry tells the truth**: We know Jerry lies. Shalonda claims Jerry tells the truth, which is incorrect. Therefore, Shalonda = **False**.} \\
\\ 
\texttt{The question asks: Does Shalonda tell the truth?} \\
\texttt{The answer is **No**.} \\
\\ 
\texttt{A: No} \\
            }
        }
    }
    & & \\ \\

    \theutterance \stepcounter{utterance}  
    & & & \multicolumn{2}{p{0.3\linewidth}}{
        \cellcolor[rgb]{0.9,0.9,0.9}{
            \makecell[{{p{\linewidth}}}]{
                \texttt{\tiny{[GM$|$GM]}}
                \texttt{Target: No} \\
            }
        }
    }
    & & \\ \\

    \theutterance \stepcounter{utterance}  
    & & & \multicolumn{2}{p{0.3\linewidth}}{
        \cellcolor[rgb]{0.9,0.9,0.9}{
            \makecell[{{p{\linewidth}}}]{
                \texttt{\tiny{[GM$|$GM]}}
                \texttt{game\_result = LOSE} \\
            }
        }
    }
    & & \\ \\

    \theutterance \stepcounter{utterance}  
    & & \multicolumn{4}{p{0.6\linewidth}}{
        \cellcolor[rgb]{0.95,0.95,0.95}{
            \makecell[{{p{\linewidth}}}]{
                \texttt{\tiny{[P1$\leftrightarrow$P]}}
                \texttt{Post{-}episode Reflexion (not sent to GM):} \\
\texttt{New plan: My previous attempt failed because I included a detailed step{-}by{-}step explanation and intermediate reasoning, whereas the target state (based on the provided examples) requires only the final answer ("Yes" or "No") immediately following the "A:" prompt. The examples show a strict pattern where the model outputs the answer directly without any preamble or logical derivation. Next time, I will skip the reasoning process entirely and output only the single word "Yes" or "No" that corresponds to the final truth value of the person in question.} \\
            }
        }
    }
    & \\ \\

\end{supertabular}
}

\end{document}
